\section{Conclusion}

This review brings together the principal frameworks, regulatory expectations, and recognition
mechanisms shaping UK higher education delivery, with a particular focus on transnational
education (TNE) involving Sri Lankan partners. It synthesises how qualification levels and
award structures align across the FHEQ and SCQF, how oversight is organised differently across
England, Scotland, Wales, and Northern Ireland, and how recognition and admissions processes
intersect with UK~ENIC comparability guidance.

For UK providers delivering through Sri Lankan partners, programme design and quality
assurance must be explicitly mapped to FHEQ level descriptors and typical credit volumes,
ensuring that progression routes and exit awards remain transparent and academically coherent.
Under the Office for Students (OfS) regulatory framework, the registered lead provider retains
full responsibility for quality and standards wherever learning is delivered; validated,
franchised, and subcontracted delivery must therefore operate within the awarding body's
regulations, monitoring cycle, student-protection arrangements, and reporting duties.

Admissions and advanced-standing decisions should rely on UK~ENIC evidence to support clear,
equitable, and defensible judgements, particularly where comparisons between SLQF and FHEQ
levels are needed. Assurance priorities for Sri Lanka include robust partner due diligence,
scheduled monitoring of academic delivery, external examining and assessment moderation,
accurate and timely handling of reportable events, and credible teach-out planning.

Overall, the review provides a structured foundation for implementing academically credible,
regulatorily compliant, and student-centred TNE provision that is intelligible to students,
regulators, and employers. It also offers a reproducible evidence base to support ongoing
governance, partnership oversight, and programme development in UK–Sri Lanka higher education
collaborations.
